% ===========================================================================
\section{Introduction and preliminaries}\label{sec:intro}
% ===========================================================================

\paragraph{Context.}
The alignment paper~\cite{icse} establishes, at design time, the chain
\[
  P \;\simul\; \VP \;\sbis_{\cons}\; \VD \;\sbis\; D,
\]
where $P$ is the physical system, $\VP$ and $\VD$ are timed-automaton views of
the Physical and Digital Twin, $\sbis_{\cons}$ is semantic alignment under the
domain knowledge $\cons=(K,\{I_P,I_D\})$, and $D$ is the deployed Digital Twin.
The paper argues that the last relation, $\VD\sbis D$, can be obtained by
construction because the DT is software under the engineer's control. This
report makes that argument concrete for one implementation: the Verified Twin
toolchain of this repository.

\paragraph{The chain proved here.}
\begin{center}
\begin{tabular}{@{}lll@{}}
\toprule
Relation & Meaning & Result \\ \midrule
$\VD \iso \IR(\VD)$ & the compiler preserves and reflects the semantics & \cref{thm:compiler} \\
$\IR(\VD) \iso \Ksem(\IR(\VD))$ & the kernel executes exactly the IR semantics & \cref{thm:kernel} \\
$\Ksem \wbis \Kprod$ & runtime bookkeeping is unobservable ($\tauact$) & \cref{thm:production} \\
$\VD \wbis \Kprod(\IR(\VD))$ & composition & \cref{thm:composition} \\
\bottomrule
\end{tabular}
\end{center}

\paragraph{Method.}
The implementation and the proof were developed together. Wherever an
implementation choice would have complicated the argument, the
implementation was simplified instead. Examples are the exact integer time
grid (\cref{sec:time}), value-type kernel configurations that map one-to-one
onto formal configurations (\cref{sec:kernel}), and a compiler that
\emph{rejects}, rather than approximates, every construct whose semantics
it cannot preserve exactly (\cref{sec:compiler}). \Cref{app:code} maps every
definition to the C++ function that implements it.

\paragraph{What is not claimed.}
Time is \emph{logical model time}. No statement is made about wall-clock
time, worst-case execution time, scheduling, network latency, clock
synchronisation, the physical plant, or the correctness of planners and
sensor models (\cref{sec:limitations}). The argument is a pen-and-paper proof:
it has not been checked by a proof assistant. Tests in the repository are
validation, not proof.

% ---------------------------------------------------------------------------
\subsection{Labelled timed transition systems}
% ---------------------------------------------------------------------------

Throughout, $\TT\subseteq\Rnn$ denotes a \emph{time domain}: a set containing
$0$ and closed under addition, ordered by $\le$. The domain used by the
implementation is fixed in \cref{sec:time}.

\begin{definition}[LTTS]\label{def:ltts}
A \emph{labelled timed transition system} over $\TT$ is a tuple
$\mathcal{T}=(S,s_0,\Sigma,\to,\Props,\lab)$ where $S$ is a set of states,
$s_0\in S$ is the initial state, $\Sigma$ is a set of discrete labels with
$\Sigma\cap\TT=\emptyset$, ${\to}\subseteq S\times(\Sigma\cup\TT)\times S$ is the
transition relation, $\Props$ is a set of atomic propositions and
$\lab:S\to2^{\Props}$ is the labelling. A designated subset
$\Sigma_\tauact\subseteq\Sigma$ contains the \emph{internal} labels.
We write $s\trans{\lambda}s'$ for $(s,\lambda,s')\in{\to}$.
\end{definition}

This is Definition~1 of the alignment paper, with the event and delay labels
written $\Sigma$ and $\TT$.

\begin{definition}[Isomorphism]\label{def:iso}
Two LTTSs $\mathcal{T}_1,\mathcal{T}_2$ over the same $\TT$ are
\emph{isomorphic}, written $\mathcal{T}_1\iso\mathcal{T}_2$, if there are
bijections $\gamma:S_1\to S_2$, $\eta:\Sigma_1\to\Sigma_2$ (with
$\eta(\Sigma_{1,\tauact})=\Sigma_{2,\tauact}$) and $\pi:\Props_1\to\Props_2$ such that
(i) $\gamma(s_{0,1})=s_{0,2}$;
(ii) $s\trans{d}s'$ iff $\gamma(s)\trans{d}\gamma(s')$ for all $d\in\TT$;
(iii) $s\trans{a}s'$ iff $\gamma(s)\trans{\eta(a)}\gamma(s')$ for all $a\in\Sigma_1$;
(iv) $p\in\lab_1(s)$ iff $\pi(p)\in\lab_2(\gamma(s))$.
\end{definition}

\begin{definition}[Weak timed bisimulation]\label{def:wtb}
Let $\mathcal{T}_1,\mathcal{T}_2$ be LTTSs over $\TT$ whose observable label
sets are identified by a bijection $\eta$ and whose propositions are
identified by $\pi$. For $\mathcal{T}_i$ write
$s\wtrans{\epsilon}s'$ if $s$ reaches $s'$ by a finite sequence of internal
transitions, $s\wtrans{a}s'$ if $s\wtrans{\epsilon}\trans{a}\wtrans{\epsilon}s'$
for an observable $a$, and $s\wtrans{d}s'$ if
$s\wtrans{\epsilon}\trans{d_1}\wtrans{\epsilon}\cdots\trans{d_k}\wtrans{\epsilon}s'$
with $d_1+\dots+d_k=d$. A relation $\mathcal{R}\subseteq S_1\times S_2$ is a
\emph{weak timed bisimulation} if $(s_{0,1},s_{0,2})\in\mathcal{R}$ and for
every $(s_1,s_2)\in\mathcal{R}$:
\begin{enumerate}[label=(W\arabic*)]
\item $p\in\lab_1(s_1)$ iff $\pi(p)\in\lab_2(s_2)$;
\item if $s_1\trans{\lambda}s_1'$ ($\lambda$ observable or a delay) there is
  $s_2\wtrans{\eta(\lambda)}s_2'$ with $(s_1',s_2')\in\mathcal{R}$, and if
  $s_1\trans{\lambda}s_1'$ with $\lambda$ internal there is
  $s_2\wtrans{\epsilon}s_2'$ with $(s_1',s_2')\in\mathcal{R}$;
\item symmetrically for transitions of $s_2$.
\end{enumerate}
$\mathcal{T}_1\wbis\mathcal{T}_2$ if such a relation exists. Taking
$\wtrans{}$ to be $\trans{}$ yields \emph{strong} timed bisimulation $\sbis$.
\end{definition}

\begin{lemma}[Standard facts]\label{lem:facts}
(a) $\mathcal{T}_1\iso\mathcal{T}_2$ implies $\mathcal{T}_1\sbis\mathcal{T}_2$
(the graph of $\gamma$ is a strong bisimulation), and strong timed bisimilarity
implies weak timed bisimilarity.
(b) $\iso$, $\sbis$ and $\wbis$ are equivalence relations; in particular they
are transitive (relational composition of bisimulations is a bisimulation).
(c) Isomorphic LTTSs satisfy the same properties expressible over their
labels and propositions.
\end{lemma}
\begin{proof}
(a) Conditions (i)--(iv) of \cref{def:iso} are exactly the strong-bisimulation
conditions for the relation $\{(s,\gamma(s))\}$, and $\trans{}\subseteq\wtrans{}$.
(b) is standard~\cite{milner}; for weak timed bisimulation the composition
argument also covers delays because $\wtrans{d}$ composes additively.
(c) follows because $\gamma$ maps runs to runs bijectively and preserves labels
and propositions.
\end{proof}
